\ej{15}{Elección de algoritmo de ordenamiento}

Usaría \textbf{Heapsort con \textsc{Heapify} iterativo}, suponiendo que los datos de una noche caben en un arreglo en memoria y que no se exige estabilidad.

\begin{enumerate}
\item \textbf{Funcionamiento.} Construye un montículo máximo sobre el arreglo. Después intercambia la raíz con la última posición del montículo, reduce su tamaño y restaura la propiedad de montículo. Como la raíz contiene el máximo pendiente, cada extracción coloca un elemento en su posición definitiva.

\item \textbf{Memoria y recursión.} Trabaja sobre el mismo arreglo y utiliza únicamente un número fijo de índices y una variable para intercambios: la memoria adicional pertenece a $\Oh{1}$. Al implementar \textsc{Heapify} mediante un ciclo, no necesita recursión ni una pila de llamadas.

\item \textbf{Tiempo garantizado.} Un montículo de $m$ elementos es un árbol binario completo de altura $\lfloor\log_2 m\rfloor$. Cada restauración desciende, como máximo, esa cantidad de niveles y realiza trabajo constante por nivel; por tanto, cuesta $\Oh{\log m}$.

La construcción realiza a lo sumo $n$ restauraciones y el ordenamiento realiza $n-1$ extracciones. Como cada montículo tiene a lo sumo $n$ elementos, el tiempo total pertenece a $\Oh{n\log n}$ en el peor caso, y esta misma cota superior vale para el promedio. Esta garantía permite limitar el cómputo nocturno independientemente del orden inicial.

\item \textbf{Comparación.}
\begin{itemize}
\item \textbf{Quicksort:} con pivotes aleatorios tiene tiempo esperado $\Th{n\log n}$, pero una sucesión de particiones de tamaños $0$ y $m-1$ produce tiempo cuadrático. Su implementación recursiva usual también puede acumular $\Th{n}$ llamadas.
\item \textbf{Mergesort:} garantiza tiempo $\Th{n\log n}$, pero su implementación habitual sobre arreglos requiere memoria auxiliar $\Th{n}$ y, si es recursiva, una pila de profundidad $\Th{\log n}$.
\item \textbf{Insertion sort:} usa memoria adicional $\Th{1}$ y puede aprovechar datos casi ordenados. Sin embargo, con datos distintos en orden inverso desplaza
$\sum_{i=2}^{n}(i-1)=n(n-1)/2$ elementos, por lo que su peor caso es cuadrático.
\end{itemize}
\end{enumerate}

\textbf{Conclusión.} Heapsort iterativo combina memoria adicional constante, ausencia de recursión y tiempo $\Oh{n\log n}$ garantizado. Es adecuado para los recursos limitados del explorador; no es estable, por lo que esta elección supone que no se necesita conservar el orden original entre datos con claves iguales.
\fin
